\chapter{Conclusion and Future directions}
\label{ch:5}
This thesis demonstrates that physically informed machine learning can be effectively used to model complex phenomena in functional materials, substantially improving the predictive capability and efficiency of conventional first-principles methods. By integrating solid physical theories with machine learning methods, we have created sophisticated surrogate models that maintain accuracy while significantly decreasing computational demands.

Our development of higher-derivative Gaussian process models allowed us to accurately reconstruct harmonic and cubic anharmonic force constants directly from DFT-calculated energies and forces. These Gaussian process force constants (GPFCs) can then be used to perform phonon calculations, such as dispersion, lifetimes, and thermal conductivity, considerably surpassing standard finite-displacement techniques in computational efficiency. Nevertheless, the current implementation encountered practical limitations, notably high memory requirements for higher-order kernel derivatives, restricting the full realisation of advanced active learning strategies. Future research should explicitly address these computational constraints by exploring improved computational techniques and alternative implementations. Moreover, extending our methodology to complex-valued phonon descriptors beyond the $\Gamma$-point, and incorporating the Atomic Cluster Expansion (ACE) framework for accelerated data generation represent promising paths forward.

In parallel, our adaptation of the ACE framework for modelling charge transfer integrals in organic molecular dimers has shown strong potential. The ACE-based method, utilising systematically constructed symmetry-adapted descriptors, showed its ability to precisely model the complex relationship between electronic couplings, molecular orientation, and separation. Challenges still remain—especially regarding model performance when confronted with sharply peaked or discontinuous transfer integral surfaces characteristic of particular molecular systems. Addressing these challenges requires additional studies aimed at enhancing radial sampling techniques for obtaining improved descriptor resolution and distribution. Altering prediction targets, such as using log-scaled values, may result in more uniform distributions and enhanced model accuracy. Moreover, incorporating ACE descriptors with sophisticated neural architectures, such as graph neural networks or message-passing frameworks \cite{https://doi.org/10.48550/arxiv.2206.07697, Batatia2025}, offers a successful technique for capturing complicated anisotropic and long-range interactions that characterise real-world systems.

Ultimately, applying these enhanced models to large-scale and realistic simulations of device-level phenomena is a crucial next step, validating their practical utility and opening new avenues for the discovery and optimisation of novel functional materials.
